% geometry-of-work.tex — Manuscript skeleton
% Geometry of Work: Shape-Guided Search over Failed and Partial Solution Attempts
%
% Target: arXiv submission (cs.AI / cs.LG / math.HO cross-list)
% Format: single-column, ~15–20 pages main text + appendices
%
\documentclass[11pt]{article}

% ── Packages ──────────────────────────────────────────────────────────
\usepackage[utf8]{inputenc}
\usepackage[T1]{fontenc}
\usepackage{amsmath,amssymb,amsthm}
\usepackage{mathtools}
\usepackage{algorithm}
\usepackage{algpseudocode}
\usepackage{booktabs}
\usepackage{graphicx}
\usepackage{hyperref}
\usepackage{cleveref}
\usepackage{xcolor}
\usepackage{enumitem}
\usepackage{natbib}
\usepackage[margin=1in]{geometry}
\usepackage{caption}
\usepackage{subcaption}

% ── Theorem environments ──────────────────────────────────────────────
\newtheorem{definition}{Definition}[section]
\newtheorem{claim}{Claim}[section]
\newtheorem{remark}{Remark}[section]

% ── Draft annotations ────────────────────────────────────────────────
\newcommand{\todo}[1]{\textcolor{red}{\textbf{[TODO: #1]}}}
\newcommand{\claimnote}[1]{\textcolor{blue}{\textbf{[CLAIM: #1]}}}
\newcommand{\evidence}[1]{\textcolor{teal}{\textbf{[EVIDENCE: #1]}}}

% ── Title ─────────────────────────────────────────────────────────────
\title{%
  Geometry of Work:\\
  Shape-Guided Search over Failed and Partial Solution Attempts%
}

\author{%
  \todo{Author names and affiliations}
}

\date{\today}

% ══════════════════════════════════════════════════════════════════════
\begin{document}
\maketitle

% ── Abstract ──────────────────────────────────────────────────────────
\begin{abstract}
When repeated work on a hard problem stops producing information, we propose
treating the work itself as data. \emph{Geometry of Work} (GoW) represents a
history of problem-solving attempts as a structured search space---a
\emph{shape-space}---whose regions, boundaries, and conserved properties contain
information about how search is failing or progressing. The methodology
separates structural shape from outcome and epistemic status, searches by
inferring and traversing regime boundaries in shape-space, and projects proposed
structural moves back into the original domain for verification. We formalize
the GoW loop (\textsc{Map}~$\to$~\textsc{Boundary}~$\to$~\textsc{Move}~$\to$
\textsc{Project}~$\to$~\textsc{Measure}), describe \texttt{newf}---an
instrumented realization with typed representations, immutable provenance, and
adversarial challenge---and report exploratory experiments on an Erd\H{o}s--Straus
study corpus. In curated-feature trials, shape-guided context produced proposals
judged closer to a withheld structural move. In discovery trials, large language
models proposed structural regularities absent from the reference ledger, though
the predeclared aggregate endpoint was not met. One model-generated hypothesis
underwent a full challenge lifecycle, surviving with a resolution-gaining
boundary refinement. The experiments do not establish mathematical validity or
general discovery effectiveness; they demonstrate the framework's epistemic
behavior and constituent capabilities.

\todo{Revise abstract after all sections are drafted.}
\end{abstract}

% ── 1. Introduction ──────────────────────────────────────────────────
\section{Introduction}
\label{sec:intro}

\subsection{The problem with repeated hard work}
\todo{Motivate: when linear search exhausts its information yield, what do
you do with the accumulated history of attempts?}

\subsection{Geometry of Work}
\todo{Core thesis in one paragraph:}

\begin{quote}
When repeated work stops producing information, treat the work itself as data.
\end{quote}

Given work items~$W$, structural descriptions~$S(W)$, and observed
outcomes~$O(W)$, reconstruct a geometry over work, identify boundaries between
outcome regimes, propose a structural move~$\Delta S$, and project that move
back into the domain:
\[
W
\xrightarrow{S}
\mathcal{G}(W)
\xrightarrow{B}
\Delta S
\xrightarrow{\Pi}
W'
\xrightarrow{O}
\mathcal{G}'(W \cup \{W'\})
\]

\subsection{Contributions}

We make four claims, deliberately ordered from broadest to narrowest:

\begin{enumerate}[label=\textbf{C\arabic*.}]
  \item \textbf{Conceptual.} A history of problem-solving work can itself be
    represented as a structured search space whose geometry contains useful
    information about how search is failing or progressing.
  \item \textbf{Methodological.} Conserved shapes, outcome regimes, boundaries,
    and structural deltas provide a disciplined way to choose informative next
    work rather than continuing linear search.
  \item \textbf{Systems.} \texttt{newf} operationalizes this methodology using
    typed, provenance-carrying representations, epistemic state tracking,
    adversarial challenge, and domain projection.
  \item \textbf{Empirical (narrow).} In project-authored Erd\H{o}s--Straus study
    corpora, LLMs can propose structural regularities not explicitly supplied as
    reference properties, and shape-guided context can produce domain proposals
    judged closer to a withheld structural move; the experiments do \emph{not}
    establish mathematical validity, general discovery effectiveness, or
    autonomous theorem solving.
\end{enumerate}

% ── 2. Motivating Example ────────────────────────────────────────────
\section{Motivating Example}
\label{sec:example}

\todo{Concrete before/after: conventional search vs.\ shape-guided search on a
small worked example. Use Erd\H{o}s--Straus or a simpler toy problem.}

% ── 3. Semantic Model ────────────────────────────────────────────────
\section{Semantic Model: The Geometry of Work}
\label{sec:semantic}

\begin{definition}[Work item]
A \emph{work item} is a triple $W_i = (S_i, O_i, E_i)$ where $S_i$ is a
structural description (\emph{shape}), $O_i \in \mathcal{O}$ is an observed
outcome, and $E_i$ is the evidence supporting the description and outcome.
\end{definition}

\begin{definition}[Structural shape]
A \emph{structural shape} $S$ is an outcome-blind description of a piece of
work: its representation, assumptions, operators, and preserved properties.
Two attempts with the same shape are the same point in shape-space even if
written in different notation; two attempts with different shapes are different
points even if they use the same words.
\end{definition}

\begin{definition}[Geometry of Work]
Given a population of work items $\{W_1, \ldots, W_n\}$, the \emph{Geometry
of Work} is the tuple
\[
\mathcal{G}(W) = (R, B, H, V, T)
\]
where:
\begin{itemize}[nosep]
  \item $R$: regimes---regions of shape-space associated with an outcome
    distribution;
  \item $B$: boundaries---structural conditions separating regimes;
  \item $H$: shape claims---conserved or discriminating structural properties;
  \item $V$: voids---unoccupied or under-sampled structural combinations;
  \item $T$: trajectories---paths through shape-space traced by successive work.
\end{itemize}
\end{definition}

\todo{Develop regime, boundary, structural delta, projection, landing point
definitions. Keep self-contained---a reader must reconstruct GoW from this
section alone.}

% ── 4. Epistemic Model ──────────────────────────────────────────────
\section{Epistemic Model: Claims, Evidence, and Promotion}
\label{sec:epistemic}

\todo{Formalize:}

\begin{definition}[Shape claim]
A \emph{shape claim} is a tuple
$\text{ShapeClaim} = (\text{shape}, \text{conditioning}, \text{role},
\text{epistemic state})$
where the four components are orthogonal axes, not a single confidence ladder.
\end{definition}

\todo{The four separations:}
\begin{enumerate}[nosep]
  \item Shape $\neq$ outcome-conditioning
  \item Observed regularity $\neq$ obstruction
  \item Model judgment $\neq$ verification
  \item Absence $\neq$ negation
\end{enumerate}

\todo{Verification hierarchy. Epistemic state transitions. Challenge operations.
Anti-fractal discipline.}

% ── 5. Shape-Guided Search ───────────────────────────────────────────
\section{Shape-Guided Search}
\label{sec:method}

\begin{algorithm}[t]
\caption{Shape-Guided Search: \textsc{Map $\to$ Boundary $\to$ Move $\to$
  Project $\to$ Measure}}
\label{alg:gow}
\begin{algorithmic}[1]
\Require Population of prior work $W$; observations $O$; domain verifier $V$
\Ensure Updated geometry $\mathcal{G}'$ and (possibly) improved candidate $W'$
\Statex
\State \textsc{Map}: Infer structural descriptions $S(W)$; reconstruct
  $\mathcal{G}(W) = (R, B, H, V, T)$
\State \textsc{Regime}: Identify outcome-conditioned regions in $\mathcal{G}$
\State \textsc{Shape}: Propose conserved or discriminating structure $H'$
\State \textsc{Challenge}: Seek counterexamples, success contrasts,
  abstraction splits, redundancy, bias
\State \textsc{Boundary}: Identify minimal structural distinction separating
  relevant regimes
\State \textsc{Move}: Choose structural delta $\Delta S$ crossing the boundary
\State \textsc{Project}: Compile $\Delta S$ into domain-native candidate
  $W' = \Pi_{\text{domain}}(\Delta S)$
\State \textsc{Verify}: Evaluate $W'$ using domain verifier $V$; record
  landing point $(O(W'), \text{strength})$
\State \textsc{Update}: Insert $W'$, $S(W')$, $O(W')$ into $\mathcal{G}$;
  mutate search policy
\State \Return to Step~1 with updated $\mathcal{G}'$
\end{algorithmic}
\end{algorithm}

\todo{Prose walkthrough of each step. Recursive challenge refinement.
Why failed crossings are informative. Failure invariants as one instrument
among several.}

% ── 6. System ────────────────────────────────────────────────────────
\section{System: \texttt{newf}}
\label{sec:system}

\todo{Architecture figure. 20--25\% of paper, not 60\%. Frame the systems
contribution as: how do you build software that lets fuzzy model-native
operations coexist with deterministic evidence and immutable provenance?}

\todo{Layers:}
\begin{enumerate}[nosep]
  \item Work ingestion (immutable, content-addressed)
  \item Structural normalization (model-assisted, epistemic status tracked)
  \item Shape-space construction (deterministic canonicalization + clustering)
  \item Shape proposal (model-native: compress, contrast, abstract)
  \item Challenge lifecycle (adversarial, code-verified)
  \item Boundary / frontier generation
  \item Domain projection
  \item Assessment (strongest-decisive-first verifier routing)
  \item Map update (re-entry + policy mutation)
\end{enumerate}

% ── 7. Experimental Method ───────────────────────────────────────────
\section{Experimental Method}
\label{sec:method-exp}

\todo{Corpus description (Erd\H{o}s--Straus synthetic blinded benchmark).
Train/target split. Assessment protocol. Provenance guarantees.}

% ── 8. Results ───────────────────────────────────────────────────────
\section{Results}
\label{sec:results}

\subsection{Experiment A: Negative controls}
\label{sec:results-negative}
\todo{System refuses to manufacture structure when no eligible failure
population exists, vocabulary resolution is inadequate, or recovery is
epistemically undecidable. Purpose: demonstrate epistemic behavior.}

\subsection{Experiment B: Curated-feature guided search (Pilot-003)}
\label{sec:results-pilot003}
\todo{Curated-feature structural hypotheses supplied and challenged.
B0 (baseline) vs B3 (structure-informed). Automatic assessment inconclusive.
Separate model review: B3 matched intended structural move. Same-family caveat.
Criterion recovery, not mathematical validity.}

\subsection{Experiment C: Shape discovery (Pilot-004)}
\label{sec:results-pilot004}
\todo{23 proposal occurrences. Reference matches. Provisionally defensible
novel occurrences. Repeated themes. Quorum disagreement. Predeclared aggregate
criterion not met. Capability finding: unencoded shape generation occurred.}

The experiment was negative under its predeclared effectiveness endpoint while
positive on a constituent capability observation: models generated plausible,
previously unencoded shape hypotheses.

\subsection{Experiment D: Novel-hypothesis challenge}
\label{sec:results-challenge}
\todo{One reference-absent Pilot-004 hypothesis (N2a or similar) through full
challenge lifecycle: ground $\to$ challenge $\to$ \{falsify, weaken, split,
survive\}. Report outcome regardless of direction.}

% ── 9. Discussion ────────────────────────────────────────────────────
\section{Discussion}
\label{sec:discussion}

\todo{What the evidence does and does not establish. Micro-failures and
recursive geometry. Human use of GoW. The interesting distinction: experiment
negative on method reliability, positive on constituent capability.}

\subsection{Complement geometry: the dual reading}
\label{sec:complement}

The forward geometry developed in \cref{sec:semantic} reads shape from the
structure of completed work---what is \emph{present} in shape-space. A dual
reading is available: infer shape from the regions of possibility-space that
accumulated work \emph{excludes}.

Let $\Omega$ denote an ambient possibility-space of structurally plausible
work on a problem. Each work item $W_i$ imposes a constraint $C_i$ on this
space---a mechanism shown not to work, a scope demonstrated to be insufficient,
a combination established as unproductive. The surviving space after $n$ pieces
of work is:
\[
\Omega_W \;=\; \Omega \;\setminus\; \bigcup_{i=1}^{n} F_i
\]
where each $F_i$ is the region excluded by the evidence from $W_i$. As work
accumulates, the residual region $\Omega_W$ may acquire a non-trivial
shape---determined not by any single attempt but by the \emph{collective
constraints} of all attempts.

We call this the \textbf{anti-vacuum}: a region that is unoccupied by prior work
but structurally determined by everything that surrounds it.

\begin{remark}[Complementarity Principle]
Work informs search in two ways: through the structures instantiated by completed
attempts (\emph{forward geometry}) and through the regions of possibility-space
those attempts exclude or constrain (\emph{complement geometry}). Sufficiently
structured negative evidence can induce a residual geometry whose conserved
properties specify candidate directions not instantiated by prior Work.
\end{remark}

The forward and complement readings yield different objects:

\begin{center}
\begin{tabular}{lll}
\toprule
& \textbf{Forward geometry} & \textbf{Complement geometry} \\
\midrule
Input & occupied work & excluded possibility-space \\
Yields & clusters, invariants, regimes & forbidden regions, cavities, corridors \\
Key object & conserved shape & residual structural requirements \\
Generator prompt & ``cross this boundary'' & ``occupy this constrained void'' \\
\bottomrule
\end{tabular}
\end{center}

The complement view reinterprets the three outcome types as different kinds of
geometric force:
\begin{itemize}[nosep]
  \item \textbf{Failure} imposes \emph{exclusion}---structural constraints on
    what remains possible.
  \item \textbf{Partial success} imposes \emph{gradient}---directional evidence
    about where improvement lies within the surviving region.
  \item \textbf{Success} imposes \emph{attractor evidence}---direct constraint
    on what structural features a solution possesses.
\end{itemize}

Under this reading, a heavily worked problem with many diverse, well-characterized
failures is not hopeless but \emph{highly constrained}: $\dim(\Omega_n) \ll
\dim(\Omega)$. Community resistance becomes an asset rather than a discouragement,
because the residual cavity may be increasingly well defined.

This gives the three modes of failed work:
\begin{quote}
Conventional search treats failed work as \emph{consumed effort}.\\
Geometry of Work treats failed work as \emph{material}.\\
Complement geometry treats failed work as \emph{pressure}.
\end{quote}

A formal \emph{constraint-collapse operator} $\mathcal{C}(\Omega \mid W) \to R$
that computes the residual region and extracts its structural properties is a
natural next primitive, but remains a hypothesis: we have not implemented or
tested it. The complement view is stated here as a theoretically motivated
extension, not as a demonstrated capability.

\todo{Expand with connection to existing void concept in forward geometry.
Note that rigorous definition of $\Omega$ is itself a research problem.}

% ── 10. Related Work ─────────────────────────────────────────────────
\section{Related Work}
\label{sec:related}

\todo{Serious comparison against: metaheuristics, novelty search /
quality-diversity, case-based reasoning, automated scientific discovery,
automated theorem proving / conjecture generation, program synthesis / CEGAR,
meta-learning, failure analysis, design-space exploration, LLM self-reflection /
debate / verifier architectures. See docs/research/related-work-matrix.md.}

% ── 11. Limitations and Threats to Validity ──────────────────────────
\section{Limitations and Threats to Validity}
\label{sec:limitations}

\todo{Single domain. LLM-dependent structural descriptions. Same-model-family
quorum. Curated corpus. No statistical power claims. Projection bottleneck.
No convergence guarantee.}

% ── 12. Future Work ──────────────────────────────────────────────────
\section{Future Work}
\label{sec:future}

\todo{Second domain. Statistical comparison against undirected search.
Formal convergence analysis. Human-in-the-loop GoW. Micro-failure geometry.
Search-policy mutation evaluation. Cross-model verification.}

% ── 13. Conclusion ───────────────────────────────────────────────────
\section{Conclusion}
\label{sec:conclusion}

\todo{Restate thesis. Summarize contributions C1--C4. Restate what the
evidence does and does not establish.}

% ── Acknowledgments ──────────────────────────────────────────────────
\section*{Acknowledgments}
\todo{Acknowledge tools, models, reviewers.}

% ── Bibliography ─────────────────────────────────────────────────────
\bibliographystyle{plainnat}
\bibliography{references}

% ══════════════════════════════════════════════════════════════════════
% APPENDICES
% ══════════════════════════════════════════════════════════════════════
\appendix

\section{Formal Definitions}
\label{app:formal}
\todo{Complete formal definitions of all objects and operators.}

\section{Protocol Details}
\label{app:protocol}
\todo{Full experiment protocols for Pilots 003 and 004.}

\section{Complete Experiment Tables}
\label{app:tables}
\todo{Full result tables with all entries, dispositions, quorum votes.}

\section{Claim/Evidence Ledger}
\label{app:claims}
\todo{Every empirical sentence in the paper mapped to its frozen artifact.}

\section{Falsification Conditions}
\label{app:falsification}
\todo{Full statement of what evidence would make GoW uninteresting or wrong.
See docs/theory/00-paper-claims.md for the five conditions.}

\end{document}
